\chapter{All ternary layers: bounded original-support reductions and a
longer crossing window}
{\small\sffamily Source: \nolinkurl{repository/collatz_reconstruction/research_program/ternary_join_reduction_20260916/note.md}.\par}


16 September 2026. Additive continuation for
\nolinkurl{KokunoYumeto/collatz-workbench}, PR \#2, read at
\nolinkurl{9afcbe2dc9bace77c1ae705b1cd128da22fe8f8b}. The current
receiver is the absolute integral first-jet complex in
\nolinkurl{intrinsic_zero_firstjet_20260915}. The local
\nolinkurl{height_window_reduction_20260916} is a source of the
preceding envelope and common-future calculus. Its pending files are not
claimed to be published by this addition. The definitions and proofs
used here are supplied below; execution depends only on the published
first-jet validator, whose original hash is checked.

The main new construction has two complete families at every positive
ternary valuation, a finite catalogue of all indicated
one-step/common-future comparisons, and a whole-source retraction with a
quadratic bound on original vertex support. The first-crossing
certificate advances from 10,280 to 10,945 odd steps. None of these
results asserts that all remaining roots vanish. No logarithm estimate,
probabilistic hypothesis, or RH estimate is used.

\section{1. Original maps and complete affine cylinders}

Use the unchanged positive-odd source X and map

\SourceMath{T(n)=(3n+1)/2^{a(n)},\qquad a(n)=\nu_2(3n+1).}{}

For a word w=(a\_1,\ldots,a\_k) put A\_j=sum\_\{i\textless=j\}a\_i,
A=A\_k, U=2\^{}A, L=3\^{}k and

\SourceMath{C=\sum_{j=0}^{k-1}3^{k-1-j}2^{A_j},\qquad f_w(x)=(Lx+C)/U.}{1}

The complete source and target are

\SourceMath{\rho+2Ut\longmapsto\eta+2Lt,\qquad t\ge0,
\quad \rho\equiv(U-C)L^{-1}\pmod{2U},\quad0<\rho<2U,
\quad\eta=(L\rho+C)/U.}{2}

The parameter inverses are (n-rho)/(2U) and (y-eta)/(2L) on these actual
images. Terminal oddness in (2) implies each original prefix oddness:
reduce the full numerator modulo 2\^{}\{A\_j+1\}; its first j terms form
the prefix numerator times an odd power of three, its next term has
valuation exactly A\_j, and all later terms vanish. Consecutive prefix
equations then give the exact valuations a\_j. Positivity propagates
from the positive starting value. The same congruences work for the
negative source rho-2U; every one of its original iterates is negative,
since 3x+1\textless0 for a negative odd integer. Therefore

\SourceMath{0<\eta<2L.}{3}

This sign argument is about the original formulas on a stated additional
source; it is not an identification of positive and negative cycles.

\section{2. Two source families at every original ternary valuation}

Fix integers b\textgreater=1 and e in \{2,6\}. Define

\SourceMath{h_e=(2^e+2)/3\in\{2,22\},\qquad
 a=a(b,e)=\min\{a\ge1:2^{a+e+2b}<3^{a+b}\}.}{4}

This minimum exists with an elementary explicit bound:

\SourceMath{a(b,e)\le b+2e,\qquad a(b+1,e)-a(b,e)\in\{0,1\}.}{5}

Indeed at a=b+2e the ratio in (4) is (8/9)\^{}\{b+e\}\textless1.
Increasing b multiplies that ratio by 4/3, and then increasing a by one
multiplies it by 2/3. The product is 8/9\textless1, proving the adjacent
increment claim along with minimality.

Keep the original right word

\SourceMath{w_{b,e}=(1)^a\,(e)\,(2)^{b-1}\,(3).}{6}

Its length is a+b+1, and its exponent sum is a+e+2b+1. Put

\SourceMath{Q=3^{a+b},\qquad J=2^{e+2b},\qquad K=2^{a+e+2b}<Q.}{}

Define the EXACT larger-source set

\SourceMath{\mathcal R_{b,e}=\{n>0:\ n\equiv3\pmod4,\quad
 Jn+h_e3^b\equiv0\pmod Q\}.}{7}

This is one full positive residue class modulo 4Q. Let n\_0 be its
representative in (0,4Q), and set

\SourceMath{m_0=\frac{Kn_0+2^ah_e3^b}{Q}-1.}{}

\subsection{Theorem 1. Complete lower-source joins on every layer}

For every b,e above, all original pairs in the designated fibre are

\SourceMath{\boxed{\begin{aligned}
n(t)&=n_0+4Qt,\\
m(t)&=m_0+4Kt,\\
y(t)&=(3n_0+1)/2+6Qt,\qquad t\in\mathbb Z_{\ge0},
\end{aligned}}}{8}

and they obey the ORIGINAL forward paths

\SourceMath{n(t)\xrightarrow{1}y(t)\xleftarrow{\ w_{b,e}\ }m(t),\qquad 0<m(t)<n(t).}{9}

Every n(t) has actual valuation nu\_3(n(t))=b. Its ternary unit
n(t)/3\^{}b is 1 modulo 3 for e=2 and 2 modulo 3 for e=6. In particular
it has no incoming original odd-return edge. The two families at
different (b,e) are disjoint. No claim is made that (7) covers every
source with that valuation and unit.

\textbf{Proof with the original coordinates.} Congruence (7) first
forces 3\^{}b\textbar n.~With z=n/3\^{}b it is Jz+h\_e=0 modulo 3\^{}a.
Since a\textgreater=1, J=1 modulo3, and h\_e is a unit modulo3, z is a
unit of the asserted residue. Thus b is exactly the valuation, not just
a divisibility lower bound. The original input is n=3\^{}b z with n=3
mod4.

Set w=(Jz+h\_e)/3\^{}a. It is a positive even integer with nu\_2(w)=1: J
is divisible by 16, h\_e=2 modulo4, and division by odd 3\^{}a preserves
valuation. Define

\SourceMath{m=2^a w-1=\frac{Kn+2^ah_e3^b}{Q}-1.}{10}

After j of the first a odd returns its ORIGINAL value is

\SourceMath{x_j=3^j2^{a-j}w-1\qquad(0\le j\le a).}{11}

For j\textless a, nu\_2(x\_j+1)\textgreater=2, so the next exponent is
exactly one. At j=a,

\SourceMath{x_a=2^e(4/3)^b n+h_e-1.}{12}

The next original exponent is exactly e, since 3h\_e-2=2\^{}e and

\SourceMath{T(x_a)=1+4^b n/3^{b-1}.}{}

For the subsequent b-1 steps, the value after j of those steps is

\SourceMath{1+4^{b-j}n/3^{b-1-j}\qquad(0\le j\le b-1).}{13}

They have exact exponent two and end at 1+4n. The last exponent is
exactly three: 3(1+4n)+1=4(3n+1), and n=3 modulo4 gives nu\_2(3n+1)=1.
The endpoint is (3n+1)/2. Thus every displayed edge is integral, odd at
its odd-return endpoints, positive, and has its original prescribed
valuation. Conversely equality of these two word endpoints, using (1),
gives (10) and hence (7). The first n=3 modulo4 condition is necessary
for the left word (1). This proves the full source domain.

For strict height, (10) reads

\SourceMath{m=(K/Q)n+h_e(2/3)^a-1.}{}

Equation (4) implies h\_e(2/3)\^{}a \textless{} (h\_e/2\^{}e)(3/4)\^{}b
\textless{} 1. Both K/Q\textless1 and this negative constant term
therefore give m\textless n, while m=2\^{}a w-1\textgreater0 gives
positivity. Since both endpoints are odd, n-m\textgreater=2. CRT between
Q and 4 gives (8); subtracting and dividing any of the three original
progressions supplies its inverse parameter. An incoming edge to n would
give 3x+1=2\^{}c n=0 modulo3, impossible. QED.

For e=2,b=1 this is the previously retained family 1731+2916t to
1215+2048t. For e=2,b=2 it gives

\SourceMath{6219+8748t\longrightarrow9329+13122t\longleftarrow5823+8192t.}{14}

For e=6,b=2, a=12, and it gives

\SourceMath{\boxed{13704327+19131876t\longrightarrow20556491+28697814t
\longleftarrow12017663+16777216t.}}{15}

The right word in (15) is (1)\^{}12,(6),(2),(3). This is a
lower-component relation, not a purported forward descent from the left
source to the right source.

\subsection{Complete treatment of longer leading runs in this template}

For a'\textgreater=a(b,e) the same formula with a' gives a subfamily of
(7), since its congruence modulo 3\^{}\{a'+b\} implies the original
congruence modulo 3\^{}\{a+b\}. Both source morphisms remain available
and generally have different smaller endpoints. The minimal a(b,e) thus
covers the complete height-reducing template family with longer leading
runs, rather than being a sample of the a-axis.

The membership test at an arbitrary n is finite: compute its ACTUAL
b=nu\_3(n), read its ternary unit, select e=2 or 6, compute a by integer
comparison, and test (7). The accepted path has at most 2b+2e+1 odd
returns on its right leg. This is a bound on the constructed witness,
not on an arbitrary orbit's stopping time.

\section{3. A complete finite catalogue of further original comparisons}

There is also a useful all-length rule for any right word w of length
k\textgreater=2 with

\SourceMath{a_k\text{ odd},\qquad 2^{A-1}<3^{k-1}.}{16}

Its target eta is 5 modulo6, because eta is odd and
eta=2\^{}\{-a\_k\}=-1 modulo3. Thus its entire image is inside the
original image of the left word (1). Using (2), every pair is

\SourceMath{\boxed{\begin{aligned}
n(t)&=(2\eta-1)/3+4\cdot3^{k-1}t,\\
m(t)&=\rho+2^{A+1}t,\\
y(t)&=\eta+2\cdot3^kt,\qquad t\ge0.
\end{aligned}}}{17}

The common endpoint, either source, and their affine lattice steps have
the same inverse parameter. These pairs have m(t)\textless n(t)
throughout. To prove it, note C\textgreater=3\^{}k-2\^{}k, by termwise
comparison A\_j\textgreater=j. For k\textgreater=2,
3\^{}k-2\^{}k\textgreater3\^{}\{k-1\}\textgreater2\^{}\{A-1\}. Define
D\_0=(C-2\^{}\{A-1\})/3\textgreater0; integrality follows from the odd
final exponent. The original equation is

\SourceMath{m=(2^{A-1}n-D_0)/3^{k-1}<n.}{18}

Its numerator is positive on (17), since m starts at the positive rho.
Equation (3) also proves 0\textless(2eta-1)/3\textless4*3\^{}\{k-1\}.
Hence this is an ENTIRE residue class without omitted small positive
representatives. No bounded search is needed for these statements at
arbitrary k.

The finite catalogue takes exactly all (16) with
2\textless=k\textless=12. For each k the exponent sum has the precise
bound A\textless=(3\^{}\{k-1\}).bit\_length(). There are 20,397 words.
Their independent composition count is

\SourceMath{\sum_{k=2}^{12}\sum_{A=k}^{\operatorname{bitlength}(3^{k-1})}
\sum_{\substack{j\ge1\text{ odd}\\j\le A-k+1}}
\binom{A-j-1}{k-2}.}{19}

Every original candidate and its active/redundant face label is
regenerated and checked. Both original affine coefficient columns are
divided by their prescribed 2-powers, retaining the odd constant and
even slope at EVERY original intermediate point. This verifies the whole
variable-source family, not merely its anchor.

\section{4. Exact residue cover, without deleting redundant faces}

The old retained retraction first uses a forward one-return descent,
then inverses on the actual images of (1), (1,2), (1,1,1,2,1,1,4), and
then the e=2,b=1 join. Its nonbase roots are exactly

\SourceMath{n\bmod36\in\{3,7,15,19,27\},\quad
n\not\equiv8731\pmod{8748},\quad n\not\equiv1731\pmod{2916}.}{20}

These follow directly by solving the original expanding affine image
congruences; the inverse of (1,2) is (8n-5)/9 and the longer inverse is
(2048n-2363)/2187. All their least admitted remaining sources are in the
positive inverse domain. The point 1 remains the separate original base
root.

At the common period P=4*3\^{}11=708588, (20) has 98,091 nonbase root
residues. Visit (17) in the declared order (period, larger-source
anchor, exponent sum, word). A family is selected exactly when it covers
at least one still-retained old root residue. Its original full family
remains its domain; prioritization restricts execution, not its
mathematical source. This yields 370 selected families and 93,025
remaining residues. The 5,066 removed residues equal the sum of the
disjoint newly covered counts. A second sweep using the UNION OF ALL
20,397 candidates independently reproduces exactly the same remaining
set. The signed row/word and all mask hashes are recorded in
verification.json's full reconstructed result. These are counts for a
specified arithmetic reduction; they are not probabilities or counts of
counterexamples.

Let R\_12 denote that explicitly regenerated residue set. The fixed new
retraction uses all the old rules, then the selected catalogue, then
Theorem 1 on its remaining inputs. Its nonbase roots are exactly

\SourceMath{\{n>1:n\bmod708588\in R_{12}\}\setminus
\bigcup_{b\ge1,e\in\{2,6\}}\mathcal R_{b,e}.}{21}

The infinite layer subtraction is an effective single test on the actual
valuation of n.~It is not approximated by imposing a finite b cap.

The old/catalogue rules already cover the whole e=2 families b=1,2,3,4
and the e=6 family b=1. Every e=2 family b\textgreater=5 and e=6 family
b\textgreater=2 is new relative to the catalogue. For b\textless=10
these assertions are finite original residue comparisons. For
b\textgreater=11, every member has 3\^{}11\textbar n and n=3 modulo4,
hence has the single residue 177147 modulo 708588; that residue lies in
R\_12. This proves the comparison at every larger b, not by
extrapolation from testing 64 layers. It also proves that no fixed
catalogue length of twelve has already supplied these infinite-layer
exclusions.

For example the catalogue gives 91+324t -\textgreater{} 71+256t as a
height relation, from the original right word (1,1,2,2,1) meeting the
left word (1). Cancelling their actual common last edge recovers the
four-edge inverse path; its cancellation is recorded, not interpreted as
absence of the source.

\section{5. Integral first-jet maps and nested retractions}

On the absolute original graph use D=Z{[}epsilon{]}/epsilon\^{}2,
q=1+epsilon, and

\SourceMath{d_\epsilon e_n=V_n-qV_{T(n)},\qquad q^j=1+j\epsilon\quad(j\in\mathbb Z).}{}

The inverse q\^{}\{-1\}=1-epsilon is integral. For original forward legs
l:n-\textgreater y and r:m-\textgreater y of lengths s,t, their weighted
chains give

\SourceMath{\boxed{d_\epsilon(B_l(n)-q^{s-t}B_r(m))=V_n-q^{s-t}V_m.}}{22}

Both legs, not only their difference, are retained by the checker.
Theorem 1 has s-t=-(a+b). The coefficient kernel receiver sends V\_n to
{[}epsilon V\_n{]}. Multiplying (22) by epsilon gives the exact
lower-height relation xi\_n=xi\_m. It does not assert that either class
is zero before a lower-source boundary is supplied.

Each chosen move has m\textless n.~Recurse on the original positive
integer, leaving roots as H(V\_r)=0 and Q(V\_r)=V\_r. For a move with
chain B and signed clock k, define

\SourceMath{H(V_n)=B+q^kH(V_m),\quad Q(V_n)=q^kQ(V_m),\quad F=I-Hd_\epsilon.}{23}

The recursion terminates on EVERY original input, independent of
Collatz. Its intermediate original path values need not decrease. The
identities are

\SourceMath{d_\epsilon H=I-Q,\quad Q^2=Q,\quad HQ=0,
\quad d_\epsilon F=Qd_\epsilon,\quad F^2=F,\quad FH=0.}{24}

The first telescopes (22). The next two hold because Q ends at roots and
H is zero there. Substituting dH=I-Q and HQ=0 in F\^{}2 and FH proves
the remaining statements. Thus the projections F,Q and the inclusion of
their image complex give an actual chain-homotopy retraction. In edge
degree its full presentation is

\SourceMath{C^0_D/\operatorname{im}H\xrightarrow{\sim}\operatorname{im}F,\qquad[e]\mapsto Fe,}{}

with inverse Fe mapped to {[}e{]}, since ker F=im H. The projected edges
have these relations; they are not declared an independent basis.

The catalogue and all-layer rules act only after the preceding rules
fail. The old route is therefore unchanged up to its retained root.
Appending the new route there gives the correction K=H\_new-H\_old,
whose original boundary is Q\_old-Q\_new and whose chain lies in the old
image: F\_old K=K. Consequently

\SourceMath{F_new F_old=F_old F_new=F_new,\qquad
Q_new Q_old=Q_old Q_new=Q_new.}{25}

These apply both at the catalogue extension and the infinite-layer
extension. They prove preservation through specific maps, not through a
comparison of dimensions. All arithmetic choices are independent of
epsilon, so constant reduction sends q\^{}k to 1 and commutes with all
the maps. The original first-jet comparison kernel and its
cycle-index/free component submodules are preserved. In particular e\_1
is retained and F(e\_1)=e\_1; no integer cycle period is inverted.

\section{6. A quantitative bound on the whole-system retraction}

\subsection{Theorem 2. Polynomial original-support control}

For every original odd input n\textgreater=3, the construction above
uses only original odd vertices at most

\SourceMath{\boxed{64n^2+21.}}{26}

Before like edge terms are collected, its homotopy column has at most

\SourceMath{\boxed{K(n)=\frac{n-1}{2}\left(2\lfloor\log_3 n\rfloor+14\right)}}{27}

signed monomial edge terms. Its constant coefficients have total
absolute sum at most K(n); its epsilon coefficients have total absolute
sum at most K(n)\^{}2. The logarithm in (27) is equivalently the largest
b with 3\^{}b\textless=n, so all bounds are integer-computable. This is
a bound on the retraction to retained roots, not on all actual stopping
times.

\textbf{Proof.} Each selected move drops between positive odd integers
by at least two, so there are at most (n-1)/2 moves. The finite
catalogue has right-leg length at most twelve, and its left leg has
length one; old inverse paths are shorter. The all-layer right leg has
a+b+1\textless=2b+2e+1\textless=2b+13 edges, and its left leg adds one.
Its actual b=nu\_3(current source)\textless=floor(log\_3 n). Thus (27)
bounds the total terms.

On an all-layer move with larger source x\textless=n, the right leg
first rises to the explicit maximum (12), then every subsequent exponent
is at least two and the values decrease. Since
(4/3)\^{}b\textless=3\^{}b\textless=x and e\textless=6,h\_e\textless=22,
(12) is at most 64x\^{}2+21\textless=64n\^{}2+21. The left target is
(3x+1)/2, smaller than that bound. For every finite catalogue or inverse
leg, its smaller source is below x and its length is at most twelve.
Each original map is bounded by g\_1(z)=(3z+1)/2 on the same positive
input, hence all intermediate values are less than
\nolinkurl{(3/2)^12(x+1)<130(x+1)<=64x^2+21} for x\textgreater=3.
Forward descent is already below x. This proves (26) along every
concatenated leg.

All weights in (23) are signed q-powers. Their exponent is a sum of
previous signed leg lengths plus the position within the current leg,
whose absolute value is at most the total number K(n) of traversed
edges. Each term consequently has constant amplitude +/-1 and epsilon
amplitude of absolute value at most K(n). Collection cannot increase
either absolute-sum bound. QED.

The induced original-height filtration maps satisfy

\SourceMath{Q(C^1_{\le n})\subset C^1_{\le n},\qquad
H(C^1_{\le n})\subset C^0_{\le64n^2+21},
\quad F(C^0_{\le n})\subset C^0_{\le256n^2+21}.}{28}

The last uses T(x)\textless2x, so each H column in Hd(e\_x) obeys the
bound at source 2n. This does not assert that H preserves height without
an increase. Under expansion to the original elementary map,
intermediate odd-step values 3x+1 are at most 192n\^{}2+64 in an H
column; the divisions do not increase that bound.

For finite declared equation support E, include its original targets and
the finite original supports of these H columns. This enlargement
preserves unions and inclusions. Equations (24)-(28) describe its actual
comparison with larger support stages. The reconstructed supported zero
keeps its original fibre and relations. An infinite diagram is being
controlled by bounds on every original column, rather than by
identifying it with an absent scalar.

\section{7. Extending the sharp crossing-height window}

Retain b\_j=bitlength(3\^{}j)-1, b\_0=0, and the exact original
coefficients

\SourceMath{C_0^*=0,\quad C_m^*=3C_{m-1}^*+2^{b_{m-1}},\quad
D_m^*=2^{b_m+1}-3^m,\quad \Theta_m=C_m^*/D_m^*.}{29}

For a first odd-return coefficient crossing at m, every proper
A\_j\textless=b\_j and the final A\_m\textgreater=b\_m+1. Thus

\SourceMath{C_m^*-C_p=\sum_{j<m}3^{m-1-j}(2^{b_j}-2^{A_j})\ge0,
\quad C_p/(2^{A_m}-3^m)\le\Theta_m.}{30}

The maximum is attained at the original word with consecutive prefix
sums b\_j and final sum b\_m+1. The differences and prefix sums are
inverse operations. This is a sharp bound on the word domain; attainment
is not an assertion of an integer non-descending source in that word's
cylinder. The original non-descent equation forces
n\textless=floor(Theta\_m).

For the shortened map S(n)=(3n+1)/2 on odd n and n/2 on even n, each
original odd edge of exponent a expands to 1 followed by a-1 even
divisions. The first shortened coefficient crossing with m odd steps
occurs at b\_m+1 total shortened steps. Its numerator remains C\_p and
its denominator is 2\^{}\{b\_m+1\}; (30) therefore applies at that
possibly even endpoint as well. This supplies the exact clock
comparison.

Define H\_M=max\_\{m\textless=M\}floor(Theta\_m). All exceptions by
horizon M lie among original odd n\textless=H\_M. This proof is valid
for every chosen M without a logarithm theorem. The
coefficient-stopping-time question itself is the established Terras
problem; we do not claim it originated here {[}RT{]}.

The previous retained window had M=10280,H\_M=5624777. The current exact
recurrence gives

\SourceMath{\boxed{H_{10945}=6728242,\quad\text{largest odd candidate }6728241.}}{31}

The attaining record is \nolinkurl{(m,b_m+1,H)=(10281,16295,6728242);}
the next record is (10946,17349,8129266). Thus testing the next original
source interval extends the entire horizon through 10945, not just the
single next value 10281.

Every original odd source 3\textless=n\textless=6728241 has been
checked, with ledger row {[}n, actual odd endpoint, complete
first-crossing exponent word{]}. There are 3,364,120 rows, 551,732 new
rows beyond the previous interval. Both the odd endpoint and the earlier
shortened crossing endpoint are strictly below n.~The latter is computed
as terminal*2\^{}\{A-bitlength(3\^{}m)\}, not inferred from its later
odd value.

The independent auditor imports neither join\_control nor the first-jet
code. It replays every individual division, every proper coefficient
prefix, and the original consecutive source labels. It checks 11,738,935
original odd-return edges and 23,482,682 divisions. The longest first
crossing in this finite source interval is 141 odd returns; this
measured statistic does not replace the proved horizon in (31).
Inductive endpoint heights give a maximum total 248 odd returns, at
source 6649279, on this certified finite interval.

It follows from (30)-(31) and the complete ledger that every original
n\textgreater=2 whose first shortened coefficient crossing uses at most
10945 odd steps descends at that crossing. Starting integers and the
final valuations are unbounded. Even starting sources descend on their
first division. The original fixed loop n=1 is retained separately. No
source with no crossing is accepted by this theorem.

Each row ends at a smaller positive odd integer. Induction in the
original source order supplies actual finite arrival certificates
throughout the interval, so the original kernel-height subgroup
generated by xi\_n={[}epsilon V\_n{]}, n\textless=6728241, is zero with
its source rows retained.

\section{8. Witness extraction and remaining global task}

The source first-jet theorem states that a finite integral witness (b,c)
obeying

\SourceMath{db=0,\qquad dc-Pb=V_n}{32}

extracts an actual path from n to 1. Its proof uses the ORIGINAL
integral component relation -km=1 and thus forces the actual cycle
length m=1. It does not allow rational division by a longer period.

From (24), multiply an H column by epsilon. It is a relation epsilon
V\_n minus epsilon V\_root. Add its constant edge coefficients to the
linear part of the root's original witness (32); leave the constant
cycle untouched. The full result is again (32) at n.~The unchanged
authenticated validator performs the extraction. Six examples, including
n=13704327 outside (31), are delivered. Its new arithmetic root is
5341183 inside (31), its signed reduction clock is -16, and its
extracted ACTUAL forward path to 1 has 89 odd returns. The root is not
misreported as an actual forward iterate at that signed clock. Infinite
progressions are proved for lower-source relations; these six point
arrival witnesses do not prove arrival of all their variable endpoints.

The complete first-jet component kernel is preserved by (24)-(25). A
nonbase cycle retains its Z/mZ component and a nonperiodic component its
Z. Its least original integer must now exceed 6728241, belong to (21),
and have no first coefficient crossing with at most 10945 odd steps.
These are necessary original-source restrictions; the remaining root set
is not claimed empty.

The next arithmetic work can enlarge either leg of the common-future
family, construct additional coefficient templates on (21), or sharpen
(30) on remaining source residues. A family must retain its exact source
domain and both inverse parameters. The polynomial support bound (26)
makes the present whole-system retraction quantitatively available for
such further comparisons; it does not supply global applicability of a
next root reduction. No unproved well-founded rank, zero measure
argument, or infinite intersection assumption is used.

\section{9. Verification, integration and sources}

The general algebraic and arithmetic-family results above have their
written proofs. The finite catalogue and original-source window have
COMPLETE bounded certificates. Regression checks additionally test 128
layer objects with b\textless=64, four source parameters per layer,
exact affine coefficients, old/new projection identities, negative
controls and original witness extraction. General b coverage is proved
by (4)-(13), not extrapolated from that test range. Both ordinary and
optimized outputs are compared byte for byte. A typed-cache bug caught
by a Boolean-as-integer negative control during development was repaired
before the accepted runs; those invalid inputs remain rejected tests,
not silently removed cases.

\nolinkurl{verify.py --output full.json} reconstructs all 370 selected
families, all masks, examples and regression records.
\nolinkurl{replay.py --full --predecessor} regenerates the complete
source ledger in both execution modes and audits every row
independently. The source ledger is an external deliverable and a
reproducible optional artifact; its hash is pinned in
\nolinkurl{window_certificate.json}. No large binary need be trusted
without this reconstruction. Source/current Git identities and
publication scope are recorded separately; an expected or scheduled CI
result is not called a pass.

{[}HW{]} Owner-directed local continuation, \emph{Global height
reduction: sharp crossing windows and original coalescence maps}, 16
September 2026. Supplied \nolinkurl{Collatz_Height_Window_Proof.md} and
its complete ZIP. The source-height envelope, all-pair common-future
formula and unbounded retraction definitions are prior work of this
workbench, reproduced where used. Exact input file hashes are in
SOURCE\_INTAKE.

{[}FJ{]} KokunoYumeto/collatz-workbench,
\nolinkurl{intrinsic_zero_firstjet_20260915}, at commit
\nolinkurl{9afcbe2dc9bace77c1ae705b1cd128da22fe8f8b.} Original
firstjet.py SHA-256
\nolinkurl{a53cef65109285f377e2c1344ae662c4a74574a30c1c40dfffebc9ed6cca06b0}.
\nolinkurl{https://github.com/KokunoYumeto/collatz-workbench/tree/9afcbe2dc9bace77c1ae705b1cd128da22fe8f8b/collatz_reconstruction/research_program/intrinsic_zero_firstjet_20260915}

{[}SZ{]} Owner's Split-Zero support-fibre and killed-representative
construction, traced in {[}FJ{]}. Its support fibre and original
relation quotient are not defined as chronological weighted-automata
histories. No additional Zeta numerical estimate is imported here, and
no claim about the complete current RH programme is made.

{[}RT{]} Olivier Rozier and Claude Terracol, \emph{Paradoxical behavior
in Collatz sequences}, arXiv:2502.00948, version 5, 17 May 2026; DOI
\nolinkurl{10.1016/j.disc.2026.115167.}
\nolinkurl{https://arxiv.org/abs/2502.00948} The
coefficient-stopping-time conjecture and paradoxical segments are
established literature. No published verification record or novelty
priority is claimed here.
